\documentclass{stex}

\libusepackage{naproche}
\libusepackage{copyright}
\libinput{preamble}

\begin{document}
\begin{smodule}{characters}

\importmodule{sections?first-order-semantics}

\ifinputref\else\setsectionlevel{section}\fi

\begin{sfragment}[id=sec:characters]{Characters}
  \symdecl*{letter}
  \symdecl*{uppercase}
  \symdecl*{lowercase}
  \symdecl*{digit}
  \symdecl*{alphanumeric}
  \symdecl*{symbol}
  %
  \begin{sparagraph}[style=symdoc,for={letter,uppercase,lowercase}]
    A \definame{letter} is a character from the Basic Latin Unicode block
    in the ranges \texttt{A}--\texttt{Z} (\unicode{0041}--\unicode{005A}) and
    \texttt{a}--\texttt{z} (\unicode{0061}--\unicode{007A}).
    A \sn{letter} in the former range is called an \definame{uppercase}
    \sn{letter} and one from the latter range a \definame{lowercase}
    \sn{letter}.
  \end{sparagraph}

  \begin{sparagraph}[style=symdoc,for=digit]
    A \definame{digit} is a character from the Basic Latin Unicode block
    in the range \texttt{0}--\texttt{9} (\unicode{0030}--\unicode{0039}).
  \end{sparagraph}

  \begin{sparagraph}[style=symdoc,for=alphanumeric]
    \Sns{letter} and \sns{digit} are called \definame{alphanumeric} characters.
  \end{sparagraph}
  
  \begin{sparagraph}[style=symdoc,for=symbol]
    A \definame{symbol} is a character from the Basic Latin Unicode block
    in the ranges
    \unicode{0021}--\unicode{002F},
    \unicode{003A}--\unicode{0040},
    \unicode{005B}--\unicode{0060} and
    \unicode{007B}--\unicode{007E},
    i.e. any character that is neither a control character, a space character,
    a \sn{digit} nor a \sn{letter}.
  \end{sparagraph}
\end{sfragment}

\ifinputref\else\printlicense[CcByNcSa]{Marcel Schütz (2026)}\fi

\end{smodule}
\end{document}
